
Para poder diseñar e implementar un circuito biohíbrido, se ha de tener en cuenta el tipo de neurona con la que se va a interactuar y su circuito natural. Uno de los circuitos más estudiados y usados en la comunidad de neurociencia es el ganglio estomatogástrico (\textit{STG}) de los crustáceos, en concreto su circuito pilórico (pues es un \textit{CPG} o generador central de patrones) (Figura \ref{FIG:CPGPILORIC}) en el que sus neuronas generan secuencias rítmicas robustas, incluso cuando se encuentra aislado del resto del sistema nervioso \cite{Selverston2000, Szcs2000}. En el CPG pilórico se conocen todas las neuronas y todas las conexiones. Esto facilita la interacción con elementos artificiales en circuitos biohíbridos, ya que la actividad de las neuronas vivas puede registrarse con precisión y utilizarse para establecer acoplamientos eléctricos o químicos mediante herramientas de simulación en tiempo real. Es importante recalcar que, además de la calibración temporal de los modelos, hay que realizar una calibración en amplitud de sus señales, y que esta puede ser específica para cada preparación \cite{Reyes-Sanchez2020, ReyesSnchez2023}.

Por tanto, un \textit{CPG} es un circuito que genera de forma autónoma ritmos consistentes de activaciones secuenciales de neuronas. Dichas activaciones secuenciales están generadas por la inhibición mutua de neuronas \cite{Selverston2000, ReyesSnchez2023}.

Es posible aislar el \textit{STG} de un cangrejo y mantenerlo vivo, registrando la actividad eléctrica de sus neuronas e inyectando corriente de acuerdo a un modelo de sinapsis con microelectrodos \cite{Reyes-Sanchez2020}. Mediante programas como \textit{RTXI} y protocolos como \textit{dynamic clamp}, se pueden incorporar neuronas artificiales al circuito vivo interactuando  en bucle cerrado con estas neuronas biológicas \cite{Amaducci2022, Reyes-Sanchez2020} (Figura \ref{FIG:SETUP}).

Como ha sido mencionado previamente, las interacciones entre el modelo y la neurona se pueden construir mediante modelos de sinapsis eléctricas o sinapsis químicas que simulan la corriente que se produce en este tipo de conexiones. A pesar de que se han implementado ambos plugins (Sección \ref{SEC:PLUGINS}), por simplicidad se ha usado solo la sinapsis eléctrica en los circuitos biohíbridos de este trabajo (Sección \ref{SEC:EJEMPLOSCIRCUITOSBIOHIBRIDOS}).

\begin{figure}[CPG pilorico]{FIG:CPGPILORIC}{En la figura se muestra el circuito pilórico de un cangrejo bajo el microscopio. Las células circulares son neuronas, a los lados se pueden observar dos electrodos, uno de registro y otro de estimulación, convergentes en una neurona.}
    \image{0.7\textwidth}{}{cpg-piloric}
\end{figure}

\begin{figure}[Setup de electrodos intra y extracelulares]{FIG:SETUP}{Setup del circuito pilórico \textit{STG} con los electrodos que permiten implementar un circuito biohíbrido registrando voltaje e introduciendo corriente en las neuronas.}
    \image{0.7\textwidth}{}{setup}
\end{figure}


